\chapter{Confinement, Screening, and Oblique Confinement}
\label{ch:global-confinement-screening-oblique}

\ConventionsEuclidean

Confinement is a statement about the behavior of extended gauge-invariant
observables and the spectrum of admissible external probes.  The basic probes
are line operators labelled by electric and magnetic charges compatible with
the global form of the gauge group.  The behavior of these probes depends on
which charges can be screened by dynamical fields and which magnetic or
dyonic objects condense.

\section{Line Charges and Screening}

Let the gauge group have maximal torus \(T\).  Electric Wilson charges are
weights of \(T\) allowed by the global form of the gauge group.  Magnetic
't Hooft charges are cocharacters \(U(1)\to T\), modulo the Weyl group, again
subject to the chosen global form.  A line label is written
\[
  \gamma=(\lambda_{\rm e},\lambda_{\rm m})
  \in \Gamma_{\rm e}\oplus\Gamma_{\rm m}.
\]
Mutual locality requires the Dirac pairing
\[
  \langle \gamma,\gamma'\rangle_D
  =
  \lambda_{\rm e}(\lambda'_{\rm m})
  -
  \lambda'_{\rm e}(\lambda_{\rm m})
  \in\mathbb Z
\]
for all simultaneously allowed line labels.

If the theory contains dynamical matter of electric charge \(\rho\), then a
Wilson line whose charge differs by \(\rho\) can end on a local dynamical
field.  The corresponding external charge is screened.  Thus the long-distance
line spectrum is a quotient of the ultraviolet charge lattice by charges of
dynamical particles, with magnetic charges treated analogously when magnetic
objects are dynamical in an effective description.

\section{Area and Perimeter Laws}

For a closed contour \(C\) and line label \(\gamma\), let \(L_\gamma(C)\) be
the renormalized line operator in the chosen regulator.  A confining area law
for this label means that, for large planar loops,
\[
  -\log \langle L_\gamma(C)\rangle
  =
  \sigma_\gamma\,\operatorname{Area}(C)
  +o(\operatorname{Area}(C)),
  \qquad
  \sigma_\gamma>0 .
\]
A perimeter law means
\[
  -\log \langle L_\gamma(C)\rangle
  =
  \mu_\gamma\,\operatorname{Length}(C)
  +o(\operatorname{Length}(C))
\]
after the line self-energy has been renormalized.  These are definitions of
asymptotic behavior of line correlators; they are not definitions of a local
order parameter.

\section{Strong-Coupling Lattice Area Law}

On a hypercubic Euclidean lattice with compact gauge group \(G\), the Wilson
action for a plaquette \(p\) is a class function of the plaquette holonomy
\(U_p\).  Its character expansion has the form
\[
  e^{-S_p(U_p)}
  =
  \sum_R c_R \chi_R(U_p),
\]
where \(R\) runs over irreducible representations.  At sufficiently strong
coupling the coefficients obey \(|c_R/c_{\mathbf 1}|\ll1\) for nontrivial
\(R\).  For a Wilson loop in representation \(R_0\), the group integrals over
links vanish unless the plaquette characters form a surface ending on the
loop.  The leading nonzero contribution is therefore proportional to
\[
  \left(\frac{c_{R_0}}{c_{\mathbf 1}}\right)^{A_{\min}(C)}
\]
up to representation-theoretic constants, where \(A_{\min}(C)\) is the
minimal lattice area bounded by \(C\).

\begin{proposition}[Lattice strong-coupling area mechanism]
\label{prop:lattice-strong-coupling-area-law}
For a pure compact lattice gauge theory in a convergent strong-coupling
character expansion, a Wilson loop in a nonscreened representation has an
area-law leading term.  The string tension is positive at sufficiently strong
coupling whenever the first nontrivial character coefficient is strictly
smaller in magnitude than the trivial coefficient.
\end{proposition}

\begin{proof}
Expand every plaquette Boltzmann factor in characters.  Haar integration over
each link projects onto gauge-invariant contractions.  A nonzero term with a
Wilson loop insertion must supply representation lines on plaquettes whose
boundary cancels the representation line on \(C\).  The minimal such
configuration is a surface of plaquettes spanning \(C\).  Its weight is a
product of nontrivial character coefficients, while the vacuum partition
function is normalized by the trivial coefficients.  Convergence of the
cluster expansion bounds larger surfaces by higher powers of the small
ratios, giving the stated leading area behavior.
\end{proof}

\section{Confinement, Higgs, and Oblique Confinement}

A phase can be described by the subgroup or sublattice of line charges that
have perimeter behavior.  Electric confinement means that purely electric
nonscreened probes have area law.  A Higgs phase with electric charge
condensation screens electric probes and confines appropriate magnetic
probes.  Oblique confinement means that the condensed objects are dyonic:
the perimeter-law sublattice is generated by charges
\[
  \gamma_c=(\lambda_{\rm e}^c,\lambda_{\rm m}^c)
\]
with both electric and magnetic components.  The confined probes are those
with nontrivial Dirac pairing against the condensed sublattice.

\begin{openproblem}[Continuum confinement criterion]
\label{op:continuum-confinement-criterion}
Formulate and prove a continuum QFT confinement criterion that uses
renormalized line operators, screening quotients, cluster properties, and a
mass gap, and that reproduces the lattice strong-coupling area mechanism
where a lattice regulator is available.  The criterion should also distinguish
electric, magnetic, and oblique confinement through the full line-charge
lattice rather than through a single Wilson-loop diagnostic.
\end{openproblem}
